\begin{helper}[A maximum odd balanced candidate exists]
Let $N$ be normal globally maximum data for $(T,e_0,\varphi)$ and suppose
that its height is $h(G)$.  There is a
\noderef{ScheideBalancedCandidate} $S$ for $N$ whose vertex-set cardinality
is at least that of every balanced candidate for $N$.
\end{helper}
\begin{proof}
Put $p=N.\operatorname{height}=h(G)$ and
$Q=\noderef{NormalTripleTrunk}(N)$.  By
\noderef{NormalGlobalMaximumTriple}, the global field of $N$ is a
\noderef{GlobalMaximumTriple}; apply \noderef{ScheideGraphHeightOdd} to
that field and use the assumed height equality to see that $p$ is odd.
The definition of $Q$ is the
exact prefix of $T$ through its first $p$ vertices and first $p-1$ edge
identities.  Restrict the Tashkinov property in the global field of $N$ to
this prefix.  For every retained edge position covered by
\noderef{TashkinovTree}, its old missing-color witness has a smaller vertex
index and hence is also retained.  Thus $Q$ is Tashkinov for
$(e_0,\varphi)$.

We check $Q$ against every clause of
\noderef{ScheideBalancedCandidate}.  Its Tashkinov clause was established
above.  Its exact-trunk clause holds because a tree-sequence is an initial
segment of itself in
\noderef{TreeSequenceInitialSegment}: use equality of lengths and the
identity casts in the vertex and edge clauses.  By
\noderef{TreeSequence}, the vertex map of $Q$ is injective, so
\noderef{TreeSequenceVertexSet} has cardinality
$Q.\operatorname{length}+1=p$, which proves the odd-order clause.  Every
used edge of $Q$ is the identically indexed edge of $T$, and therefore the
used-color clause holds:
$\noderef{TreeUsedColors}(\varphi,Q)\subseteq
\noderef{TreeUsedColors}(\varphi,T)$.  Finally, $Q$ has order $p$, so there
is no integer $j$ with
$p<2j<Q.\operatorname{length}+1=p$; its post-trunk pairing condition is
vacuous.  Hence the candidate family is nonempty.

For any tree-sequence $R$, injectivity of its vertex map in
\noderef{TreeSequence} identifies $|V(R)|$ with
$R.\operatorname{length}+1$ and bounds it by the finite number $|V(G)|$.
Consequently the set
\[
 \{\,|V(R)|:R\text{ is a balanced candidate for }N\,\}
\]
is a nonempty subset of the finite set
$\{0,1,\ldots,|V(G)|\}$.  It has a largest member $m$.  By the definition
of this set there is an actual balanced candidate $S$ with $|V(S)|=m$.
For every other candidate $R$, membership of $|V(R)|$ in the same set and
maximality of $m$ give $|V(R)|\le |V(S)|$, as required.
\end{proof}
